% Options for packages loaded elsewhere
% Options for packages loaded elsewhere
\PassOptionsToPackage{unicode}{hyperref}
\PassOptionsToPackage{hyphens}{url}
\PassOptionsToPackage{dvipsnames,svgnames,x11names}{xcolor}
%
\documentclass[
]{report}
\usepackage{xcolor}
\usepackage[top=30mm,left=20mm,right=20mm,bottom=30mm]{geometry}
\usepackage{amsmath,amssymb}
\setcounter{secnumdepth}{5}
\usepackage{iftex}
\ifPDFTeX
  \usepackage[T1]{fontenc}
  \usepackage[utf8]{inputenc}
  \usepackage{textcomp} % provide euro and other symbols
\else % if luatex or xetex
  \usepackage{unicode-math} % this also loads fontspec
  \defaultfontfeatures{Scale=MatchLowercase}
  \defaultfontfeatures[\rmfamily]{Ligatures=TeX,Scale=1}
\fi
\usepackage{lmodern}
\ifPDFTeX\else
  % xetex/luatex font selection
\fi
% Use upquote if available, for straight quotes in verbatim environments
\IfFileExists{upquote.sty}{\usepackage{upquote}}{}
\IfFileExists{microtype.sty}{% use microtype if available
  \usepackage[]{microtype}
  \UseMicrotypeSet[protrusion]{basicmath} % disable protrusion for tt fonts
}{}
\makeatletter
\@ifundefined{KOMAClassName}{% if non-KOMA class
  \IfFileExists{parskip.sty}{%
    \usepackage{parskip}
  }{% else
    \setlength{\parindent}{0pt}
    \setlength{\parskip}{6pt plus 2pt minus 1pt}}
}{% if KOMA class
  \KOMAoptions{parskip=half}}
\makeatother
% Make \paragraph and \subparagraph free-standing
\makeatletter
\ifx\paragraph\undefined\else
  \let\oldparagraph\paragraph
  \renewcommand{\paragraph}{
    \@ifstar
      \xxxParagraphStar
      \xxxParagraphNoStar
  }
  \newcommand{\xxxParagraphStar}[1]{\oldparagraph*{#1}\mbox{}}
  \newcommand{\xxxParagraphNoStar}[1]{\oldparagraph{#1}\mbox{}}
\fi
\ifx\subparagraph\undefined\else
  \let\oldsubparagraph\subparagraph
  \renewcommand{\subparagraph}{
    \@ifstar
      \xxxSubParagraphStar
      \xxxSubParagraphNoStar
  }
  \newcommand{\xxxSubParagraphStar}[1]{\oldsubparagraph*{#1}\mbox{}}
  \newcommand{\xxxSubParagraphNoStar}[1]{\oldsubparagraph{#1}\mbox{}}
\fi
\makeatother


\usepackage{longtable,booktabs,array}
\usepackage{calc} % for calculating minipage widths
% Correct order of tables after \paragraph or \subparagraph
\usepackage{etoolbox}
\makeatletter
\patchcmd\longtable{\par}{\if@noskipsec\mbox{}\fi\par}{}{}
\makeatother
% Allow footnotes in longtable head/foot
\IfFileExists{footnotehyper.sty}{\usepackage{footnotehyper}}{\usepackage{footnote}}
\makesavenoteenv{longtable}
\usepackage{graphicx}
\makeatletter
\newsavebox\pandoc@box
\newcommand*\pandocbounded[1]{% scales image to fit in text height/width
  \sbox\pandoc@box{#1}%
  \Gscale@div\@tempa{\textheight}{\dimexpr\ht\pandoc@box+\dp\pandoc@box\relax}%
  \Gscale@div\@tempb{\linewidth}{\wd\pandoc@box}%
  \ifdim\@tempb\p@<\@tempa\p@\let\@tempa\@tempb\fi% select the smaller of both
  \ifdim\@tempa\p@<\p@\scalebox{\@tempa}{\usebox\pandoc@box}%
  \else\usebox{\pandoc@box}%
  \fi%
}
% Set default figure placement to htbp
\def\fps@figure{htbp}
\makeatother





\setlength{\emergencystretch}{3em} % prevent overfull lines

\providecommand{\tightlist}{%
  \setlength{\itemsep}{0pt}\setlength{\parskip}{0pt}}



 


\makeatletter
\@ifpackageloaded{caption}{}{\usepackage{caption}}
\AtBeginDocument{%
\ifdefined\contentsname
  \renewcommand*\contentsname{Table of contents}
\else
  \newcommand\contentsname{Table of contents}
\fi
\ifdefined\listfigurename
  \renewcommand*\listfigurename{List of Figures}
\else
  \newcommand\listfigurename{List of Figures}
\fi
\ifdefined\listtablename
  \renewcommand*\listtablename{List of Tables}
\else
  \newcommand\listtablename{List of Tables}
\fi
\ifdefined\figurename
  \renewcommand*\figurename{Figure}
\else
  \newcommand\figurename{Figure}
\fi
\ifdefined\tablename
  \renewcommand*\tablename{Table}
\else
  \newcommand\tablename{Table}
\fi
}
\@ifpackageloaded{float}{}{\usepackage{float}}
\floatstyle{ruled}
\@ifundefined{c@chapter}{\newfloat{codelisting}{h}{lop}}{\newfloat{codelisting}{h}{lop}[chapter]}
\floatname{codelisting}{Listing}
\newcommand*\listoflistings{\listof{codelisting}{List of Listings}}
\makeatother
\makeatletter
\makeatother
\makeatletter
\@ifpackageloaded{caption}{}{\usepackage{caption}}
\@ifpackageloaded{subcaption}{}{\usepackage{subcaption}}
\makeatother
\usepackage{bookmark}
\IfFileExists{xurl.sty}{\usepackage{xurl}}{} % add URL line breaks if available
\urlstyle{same}
\hypersetup{
  colorlinks=true,
  linkcolor={blue},
  filecolor={Maroon},
  citecolor={Blue},
  urlcolor={Blue},
  pdfcreator={LaTeX via pandoc}}


\author{}
\date{}
\begin{document}

\renewcommand*\contentsname{Table of contents}
{
\hypersetup{linkcolor=}
\setcounter{tocdepth}{2}
\tableofcontents
}

\chapter{AI Capability Checkpoints}\label{ai-capability-checkpoints}

This document provides reflection scaffolds aligned to a six-domain AI
Capability model. These checkpoints are designed to support responsible,
transparent, and professional research practice in this synthetic
example.

\begin{center}\rule{0.5\linewidth}{0.5pt}\end{center}

\section{1. Awareness \& Orientation}\label{awareness-orientation}

If AI tools were used at this stage of the project:

\begin{itemize}
\tightlist
\item
  \textbf{Tools Used}: A Large Language Model (LLM) was used as a
  drafting assistant to help shape the narrative structure of the
  synthetic protocol and generate the Python code for the synthetic
  dataset.
\item
  \textbf{Limitations}: Known limitations included a tendency for the
  LLM to hallucinate overly complex or clinically specific variables
  that were not appropriate for a generic synthetic dataset.
\item
  \textbf{Checks Performed}: Human researchers carefully reviewed all
  generated covariates and code against standard epidemiological
  principles before finalising the dataset.
\end{itemize}

\begin{center}\rule{0.5\linewidth}{0.5pt}\end{center}

\section{2. Human--AI Co-Agency}\label{humanai-co-agency}

If applicable:

\begin{itemize}
\tightlist
\item
  \textbf{Tasks Supported}: Drafting boilerplate protocol text,
  commenting code, and structuring governance documents (like the risk
  register).
\item
  \textbf{Human Judgement}: The choice of logistic regression as the
  primary analytical method, and the specific distribution shapes of the
  synthetic data, remained entirely under human control.
\item
  \textbf{Human Oversight}: All AI-drafted sections were iteratively
  edited by human authors to ensure a neutral, governance-oriented tone.
\item
  \textbf{Evaluating Suggestions}: Suggestions to include complex
  machine learning algorithms (e.g., neural networks) in the methodology
  were rejected to maintain the educational simplicity of the example.
\end{itemize}

\begin{center}\rule{0.5\linewidth}{0.5pt}\end{center}

\section{3. Applied Practice \&
Innovation}\label{applied-practice-innovation}

If AI tools were used:

\begin{itemize}
\tightlist
\item
  \textbf{Specific Ways AI Supported}: AI significantly accelerated the
  creation of the \texttt{data/codebook.csv} format and helped generate
  the repetitive markdown structure of the Quarto \texttt{.qmd} files.
\item
  \textbf{Independent Elements}: The core concept (investigating missed
  appointments) and the specific relationships injected into the
  synthetic data (e.g., \texttt{previous\_missed} being strongly
  predictive) were developed independently by human researchers.
\item
  \textbf{Modifications}: AI outputs were substantially modified to
  remove any enthusiastic or promotional language regarding AI use
  itself, in line with governance requirements.
\end{itemize}

\begin{center}\rule{0.5\linewidth}{0.5pt}\end{center}

\section{4. Ethics, Equity \& Impact}\label{ethics-equity-impact}

If AI tools were involved:

\begin{itemize}
\tightlist
\item
  \textbf{Data Shared}: No sensitive or identifiable data was shared
  with AI systems; the entire project relies on simulated data.
\item
  \textbf{Bias/Fairness Risks}: When prompting the AI to draft the
  `Interpretation' section, there was a risk of the model attributing
  missed appointments solely to patient non-compliance.
\item
  \textbf{Equity Implications Review}: The research team actively guided
  the AI drafting process to ensure structural factors (like lead time
  and deprivation quintile) were highlighted, mitigating the risk of
  framing bias.
\end{itemize}

\begin{center}\rule{0.5\linewidth}{0.5pt}\end{center}

\section{5. Decision-Making \&
Governance}\label{decision-making-governance}

If AI tools were used:

\begin{itemize}
\tightlist
\item
  \textbf{Documentation}: AI use was documented transparently in the `AI
  Capability Checkpoint' callouts embedded at the bottom of each
  relevant content page.
\item
  \textbf{Disclosure Requirements}: We considered standard journal AI
  disclosure requirements and ensured this standalone document serves as
  a comprehensive record of AI assistance.
\item
  \textbf{Justifications}: The decision to use AI to generate the
  synthetic dataset rather than hand-crafting it was justified by the
  need for reproducibility and scale (N=10,000) within a demonstration
  context.
\end{itemize}

\begin{center}\rule{0.5\linewidth}{0.5pt}\end{center}

\section{6. Reflection, Learning \&
Renewal}\label{reflection-learning-renewal}

Optional reflective prompts:

\begin{itemize}
\tightlist
\item
  \textbf{What Worked Well}: Using AI to overcome the ``blank page''
  problem for governance documentation (e.g., generating the initial
  structure of the Model Card) was highly efficient.
\item
  \textbf{Future Changes}: In future iterations, we would rely less on
  open-ended conversational prompts and more on highly structured
  templates when asking the AI to synthesize data.
\item
  \textbf{Influence}: AI positively influenced the structural
  consistency of the markdown files but required careful editing to
  ensure the framing remained strictly neutral.
\end{itemize}

\begin{center}\rule{0.5\linewidth}{0.5pt}\end{center}

\section{Important Notes}\label{important-notes}

\begin{itemize}
\tightlist
\item
  AI use is \textbf{not required} for this template.
\item
  These checkpoints are designed to support transparency and governance.
\item
  Responsibility for all research outputs remains with the human
  authors.
\end{itemize}




\end{document}
